\section{What the reader should remember}

This chapter is only a mathematical preparation. It does not contain the physical theory itself. Its purpose is to make the later chapters readable.

\begin{summarybox}
The reader should remember the following points:
\begin{itemize}
    \item Scientific notation is necessary for reading very large and very small numbers.
    \item A quantity has a numerical value, a unit, and a meaning.
    \item Variables, constants, and parameters play different roles in formulas.
    \item A function describes dependence.
    \item Cartesian coordinates describe points after a coordinate system has been chosen.
    \item A graph represents a relation visually.
    \item Parametric descriptions allow us to represent curves using an additional variable.
    \item Vectors have components, magnitude, and direction.
    \item Matrices organize numbers and can represent transformations.
    \item A derivative describes local change.
    \item An integral describes accumulation.
    \item Differentiation and integration are inverse operations in the sense of calculus.
    \item A differential equation is an equation for an unknown function involving derivatives.
    \item Approximations and limiting cases help us understand expressions and check results.
\end{itemize}
\end{summarybox}

The specific meaning of these tools will be developed later. For now, the important point is that the reader can recognize the tools and understand their mathematical role.
